\section*{Introduzione}
Un corriere deve partire dal deposito, visitare tutte le consegne una volta
e tornare al punto di partenza, percorrendo la minore distanza possibile.
Questo è il problema del commesso viaggiatore, o TSP. Nel progetto i luoghi
sono punti su un piano e le distanze sono misurate in linea retta, in unità
arbitrarie: non rappresentano tempi di percorrenza su strade reali.

La domanda studiata è: \emph{una previsione appresa può aiutare a scegliere
quali percorsi approfondire, quando il tempo per cercare è limitato?}
Confronto tre modalità di ricerca:
\begin{itemize}
\item \textbf{A, ricerca di base:} esplora i percorsi e scarta quelli già
  troppo costosi per migliorare il giro migliore trovato.
\item \textbf{B, ricerca predittiva:} usa un modello allenato in precedenza
  per esplorare prima le continuazioni che sembrano meno costose.
\item \textbf{C, ricerca predittiva regolata:} aggiunge un regolatore che
  cambia il prossimo candidato quando la ricerca non migliora da molti passi.
\end{itemize}

I casi di studio sono:
\begin{itemize}
\item \textbf{Dieci mappe da cento consegne:} confronto principale dei tre
  metodi, con uno e tre secondi disponibili per ogni ricerca.
\item \textbf{Una mappa da sei consegne:} caso piccolo in cui è possibile
  verificare il giro ottimo e osservare un errore nell'ordine delle previsioni.
\item \textbf{Una mappa personale da cento consegne:} prova esplorativa
  dell'effetto di sostituire il modello appreso.
\end{itemize}

Nella prova principale C riduce mediamente la distanza rispetto a B del
$2{,}04\%$ a un secondo e del $3{,}07\%$ a tre secondi, ma peggiora alcune
mappe. Le otto sezioni seguenti distinguono il contributo della previsione
da quello della regolazione e discutono anche questi esiti sfavorevoli.
